%============================================================
% CH2 · LESSON 1 — RECURSION AND THE CALL STACK
%============================================================
\subsection{Recursion and the Call Stack}
\label{ssec:ch2-call-stack}

Chapter~1 used DFS only to answer ``which vertices can I reach?''. For that question the \emph{order} of the visit did not matter --- BFS gave the same set. This chapter is about the order. DFS has a very particular rhythm: it \textbf{enters} a vertex, disappears into the first unexplored neighbour, comes back, disappears into the next one, and finally \textbf{leaves} the vertex for good. Almost everything in this chapter --- timestamps, the DFS tree, cycle detection, subtree sizes --- is read off that rhythm of entering and leaving.

% local drawing of the running example D: a pic whose vertex / edge styles are overridden per snapshot
\tikzset{
    g2-v1/.style=gv, g2-v2/.style=gv, g2-v3/.style=gv, g2-v4/.style=gv,
    g2-v5/.style=gv, g2-v6/.style=gv, g2-v7/.style=gv, g2-v8/.style=gv,
    g2-e12/.style=gde, g2-e15/.style=gde, g2-e16/.style=gde, g2-e23/.style=gde, g2-e24/.style=gde,
    g2-e31/.style=gde, g2-e43/.style=gde, g2-e54/.style=gde, g2-e56/.style=gde, g2-e67/.style=gde,
    g2-e76/.style=gde, g2-e87/.style=gde,
    g2-D/.pic={
        \node[g2-v1] (d1) at (0,0) {1};
        \node[g2-v2] (d2) at (1.6,0.9) {2};
        \node[g2-v3] (d3) at (3.2,1.5) {3};
        \node[g2-v4] (d4) at (3.2,0.2) {4};
        \node[g2-v5] (d5) at (1.6,-1.0) {5};
        \node[g2-v6] (d6) at (3.2,-1.0) {6};
        \node[g2-v7] (d7) at (4.8,-1.0) {7};
        \node[g2-v8] (d8) at (4.8,0.2) {8};
        \draw[g2-e12] (d1) -- (d2); \draw[g2-e15] (d1) -- (d5);
        \draw[g2-e16] (d1) to[bend right=48] (d6);
        \draw[g2-e23] (d2) -- (d3); \draw[g2-e24] (d2) -- (d4);
        \draw[g2-e31] (d3) to[bend right=40] (d1);
        \draw[g2-e43] (d4) -- (d3); \draw[g2-e54] (d5) -- (d4); \draw[g2-e56] (d5) -- (d6);
        \draw[g2-e67] (d6) to[bend left=22] (d7); \draw[g2-e76] (d7) to[bend left=22] (d6);
        \draw[g2-e87] (d8) -- (d7);
    },
    % G3 of Chapter 1 (undirected), redrawn here so that this chapter compiles on its own
    g2-u1/.style=gv, g2-u2/.style=gv, g2-u3/.style=gv, g2-u4/.style=gv, g2-u5/.style=gv, g2-u6/.style=gv, g2-u7/.style=gv,
    g2-f12/.style=ge, g2-f14/.style=ge, g2-f23/.style=ge, g2-f25/.style=ge, g2-f45/.style=ge, g2-f56/.style=ge,
    g2-G3/.pic={
        \node[g2-u1] (w1) at (0,0) {1};
        \node[g2-u2] (w2) at (1.4,0.9) {2};
        \node[g2-u4] (w4) at (1.4,-0.9) {4};
        \node[g2-u3] (w3) at (2.9,0.9) {3};
        \node[g2-u5] (w5) at (2.9,-0.9) {5};
        \node[g2-u6] (w6) at (4.3,-0.9) {6};
        \node[g2-u7] (w7) at (4.3,0.9) {7};
        \draw[g2-f12] (w1) -- (w2); \draw[g2-f14] (w1) -- (w4); \draw[g2-f23] (w2) -- (w3);
        \draw[g2-f25] (w2) -- (w5); \draw[g2-f45] (w4) -- (w5); \draw[g2-f56] (w5) -- (w6);
    }
}

\subsubsection*{The running example $D$}

Part~A uses one small directed graph throughout. It is drawn so that it contains everything a DFS can meet: two cycles, a vertex that is reached twice, an edge that ``jumps ahead'' and a vertex (8) that cannot be reached from 1.

\begin{diagrambox}
\begin{tikzpicture}[x=1.15cm, y=1cm]
    \pic {g2-D};
    \node[draw=darkgray, fill=lightgray!60, rounded corners=2pt, font=\ttfamily\scriptsize, align=left, inner sep=3pt, anchor=west] (inp) at (6.6,0.1)
        {8 12\\1 2\\1 5\\1 6\\2 3\\2 4\\3 1\\4 3\\5 4\\5 6\\6 7\\7 6\\8 7};
    \node[font=\scriptsize\color{darkgray}, above=1pt of inp] {input (directed)};
    \node[note, anchor=west, text width=3.4cm, align=left] at (8.3,0.1) {adj[1] = 2 5 6\\adj[2] = 3 4\\adj[3] = 1\\adj[4] = 3\\adj[5] = 4 6\\adj[6] = 7\\adj[7] = 6\\adj[8] = 7};
\end{tikzpicture}
\figcap{The directed graph $D$: $n = 8$, $m = 12$. Neighbours are tried in the order of the adjacency lists, i.e.\ in input order; this fixes every number in Part~A. Cycles: $1 \to 2 \to 3 \to 1$ and $6 \rightleftarrows 7$.}
\end{diagrambox}

%------------------------------------------------------------
\subsubsection*{What the call stack holds}

Here is the recursive DFS of Chapter~1, with one change: it now does something both when it \emph{enters} a vertex and when it \emph{leaves} it.

\begin{lstlisting}[style=cpp]
#include <bits/stdc++.h>
using namespace std;

// Recursive DFS with the two "hooks": code before the loop runs when u is ENTERED,
// code after the loop runs when u is LEFT (all its neighbours are done).
vector<vector<int>> adj;
vector<bool> visited;
vector<int> pre, post;                 // entry order and exit order

void dfs(int u) {
    visited[u] = true;
    pre.push_back(u);                  // hook 1: pre-order (entering u)
    for (int v : adj[u])
        if (!visited[v])
            dfs(v);                    // while dfs(v) runs, u waits on the call stack
    post.push_back(u);                 // hook 2: post-order (leaving u)
}

int main() {
    int n, m;
    cin >> n >> m;
    adj.assign(n + 1, vector<int>());
    for (int i = 0; i < m; i++) {
        int a, b;
        cin >> a >> b;
        adj[a].push_back(b);           // directed edge a -> b
    }
    visited.assign(n + 1, false);
    for (int s = 1; s <= n; s++)       // restart from every vertex not reached yet
        if (!visited[s]) dfs(s);
    cout << "pre :";
    for (int v : pre) cout << " " << v;
    cout << "\npost:";
    for (int v : post) cout << " " << v;
    cout << "\n";
}
\end{lstlisting}

\begin{lstlisting}[style=trace]
input : the input of D above
output: pre : 1 2 3 4 5 6 7 8
        post: 3 4 2 7 6 5 1 8
\end{lstlisting}

Every call \texttt{dfs(u)} owns a \textbf{stack frame} that stores two things: the vertex $u$, and \emph{how far the loop over \texttt{adj[u]} has got}. When \texttt{dfs(v)} returns, the frame of $u$ resumes exactly where it stopped. The frames that are open at a given moment form the \textbf{call stack}, and the call stack always spells a path in the graph: from the start vertex, along the edges through which DFS went deeper, down to the vertex currently being explored.

\begin{diagrambox}
\begin{tikzpicture}[x=0.6cm, y=0.6cm,
    fr/.style={draw=recursionred, fill=recursionred!14, minimum width=1.25cm, minimum height=0.42cm,
               font=\scriptsize\ttfamily, inner sep=1pt},
    frtop/.style={fr, fill=currentfill, draw=accentorange, line width=0.9pt}]
    \foreach \k/\lbl/\st in {0/{enter 1}/{1}, 1/{enter 2}/{1,2}, 2/{enter 3}/{1,2,3}, 3/{leave 3, enter 4}/{1,2,4},
                             4/{leave 4, 2; enter 5}/{1,5}, 5/{enter 6}/{1,5,6}, 6/{enter 7}/{1,5,6,7}, 7/{leave 7, 6, 5, 1}/{},
                             8/{restart: enter 8}/{8}}{
        \begin{scope}[xshift=\k*1.55cm]
            \foreach \v [count=\h from 0] in \st {
                \node[fr] at (0,\h*0.75) {dfs(\v)};
            }
            \draw[darkgray, thick] (-1.15,-0.45) -- (1.15,-0.45);
            \node[font=\tiny, text width=1.45cm, align=center, anchor=north] at (0,-0.6) {\lbl};
        \end{scope}
    }
    \node[font=\scriptsize\bfseries, text=recursionred, anchor=west] at (-1.2,3.6) {call stack (top = vertex being explored)};
\end{tikzpicture}
\figcap{The call stack of the DFS on $D$ after each step. The deepest moment is \texttt{1 5 6 7}: four open frames. After vertex 1 is left the stack is empty, and the outer loop restarts the search from 8, the only vertex not reached from 1.}
\end{diagrambox}

\begin{insightbox}[The call stack is the path from the root]
At every moment the open frames, read from bottom to top, are a path $s = x_0 \to x_1 \to \dots \to x_k = u$ in the graph, where each step $x_{i} \to x_{i+1}$ is the edge through which DFS entered $x_{i+1}$. A vertex is \emph{on the stack} exactly between the moment it is entered and the moment it is left. Lessons~\ref{ssec:ch2-timestamps}--\ref{ssec:ch2-cycles} turn this one sentence into the parenthesis theorem, the edge types and cycle detection.
\end{insightbox}

\subsubsection*{Two hooks, two orders}

\begin{itemize}[leftmargin=*]
    \item \textbf{Pre-order} (code before the loop): a vertex is processed \emph{before} everything below it. Use it to pass information \emph{down}: depths, the path from the root, a value accumulated from the root (Problem~\ref{ssec:ch2-ki}).
    \item \textbf{Post-order} (code after the loop): a vertex is processed \emph{after} everything below it. Use it to collect information \emph{up}: subtree sizes, counts, sums (Pattern~2), and the finishing order behind topological sorting (Chapter~5).
\end{itemize}
On $D$ the pre-order is $1,2,3,4,5,6,7,8$ and the post-order is $3,4,2,7,6,5,1,8$. Vertex 1 is first in pre-order and almost last in post-order: everything reachable from it is entered after it and left before it. That observation, made precise, is the parenthesis theorem of the next lesson.

%------------------------------------------------------------
\subsubsection*{Iterative DFS that also knows when a vertex is left}

The iterative DFS of Chapter~1 replaced the call stack by a vector \texttt{st} and remembered, for every vertex, the index \texttt{it[u]} of the next neighbour to try --- that index \emph{is} the ``how far the loop has got'' part of a stack frame. The vertex is left when \texttt{it[u]} reaches the end of its list. So both hooks are available, and the program below prints the same two orders as the recursive one.

\begin{lstlisting}[style=cpp]
#include <bits/stdc++.h>
using namespace std;

// Iterative DFS with entry AND exit events. st[] plays the call stack, it[u] the loop position.
// Same order as the recursive version, no recursion depth limit.
int main() {
    int n, m;
    cin >> n >> m;
    vector<vector<int>> adj(n + 1);
    for (int i = 0; i < m; i++) {
        int a, b;
        cin >> a >> b;
        adj[a].push_back(b);
    }
    vector<bool> visited(n + 1, false);
    vector<size_t> it(n + 1, 0);
    vector<int> st, pre, post;
    for (int s = 1; s <= n; s++) {
        if (visited[s]) continue;
        visited[s] = true; pre.push_back(s); st.push_back(s);      // enter s
        while (!st.empty()) {
            int u = st.back();
            if (it[u] < adj[u].size()) {
                int v = adj[u][it[u]++];
                if (!visited[v]) {
                    visited[v] = true; pre.push_back(v);              // enter v ("call" dfs(v))
                    st.push_back(v);
                }
            } else {
                post.push_back(u);                                    // leave u ("return")
                st.pop_back();
            }
        }
    }
    cout << "pre :";
    for (int v : pre) cout << " " << v;
    cout << "\npost:";
    for (int v : post) cout << " " << v;
    cout << "\n";
}
\end{lstlisting}

\begin{lstlisting}[style=trace]
input : the input of D above
output: pre : 1 2 3 4 5 6 7 8          (identical to the recursive version)
        post: 3 4 2 7 6 5 1 8
\end{lstlisting}

This is the version kept in the book's library (\texttt{code/lib/dfs\_iterative.hpp}, Appendix~A). It costs \VE{} time like the recursive one, and \texttt{st} never holds more than one copy of a vertex, so its memory is \bigO{n}.

\subsubsection*{Two tempting shortcuts that are \emph{not} DFS}

Many tutorials write ``iterative DFS'' by pushing all neighbours at once. Both variants below are fine for plain reachability, but neither gives a real DFS.

\begin{diagrambox}
\begin{tikzpicture}[x=0.95cm, y=0.95cm]
    \begin{scope}
        \tikzset{g2-e16/.style={gtree, -{Stealth[length=2.4mm]}}, g2-e12/.style={gtree, -{Stealth[length=2.4mm]}},
                 g2-e15/.style={gtree, -{Stealth[length=2.4mm]}}, g2-e67/.style={gtree, -{Stealth[length=2.4mm]}},
                 g2-e54/.style={gtree, -{Stealth[length=2.4mm]}}, g2-e43/.style={gtree, -{Stealth[length=2.4mm]}},
                 g2-v6/.style=gvB}
        \pic {g2-D};
        \node[font=\small\bfseries, anchor=west] at (-0.4,2.45) {(a) mark when pushed};
        \node[font=\scriptsize, text width=6.3cm, align=left, anchor=north west] at (-0.4,-2.0)
            {Pop $u$, then push every unmarked neighbour and mark it at once. Order $1,2,5,6,7,4,3$ (8 not reached from 1). Vertex 6 gets parent 1 (red) although a real DFS enters it from 5: the ``tree'' is wrong.};
    \end{scope}
    \begin{scope}[xshift=7.6cm]
        \tikzset{g2-v6/.style=gvB}
        \pic {g2-D};
        \node[font=\small\bfseries, anchor=west] at (-0.4,2.45) {(b) mark when popped};
        \node[font=\scriptsize, text width=6.3cm, align=left, anchor=north west] at (-0.4,-2.0)
            {Pop $u$; skip it if marked; else mark it and push all unmarked neighbours (in reverse, so the first neighbour is on top). Order $1,2,3,4,5,6,7$ --- correct here. But vertex 6 sits on the stack \textbf{twice} (pushed by 1 and by 5), the stack is not a path, and nothing tells you when a vertex is \emph{left}.};
    \end{scope}
\end{tikzpicture}
\figcap{The two shortcuts on $D$ (orders produced by running them). Neither has an exit event, so neither can compute post-order values, exit times or cycles.}
\end{diagrambox}

\begin{warningbox}[Use the iterator version whenever the order matters]
Variant (a) is not even a DFS order. Variant (b) visits vertices in a DFS order, but its stack can hold up to $m$ entries and it cannot tell when a vertex is finished. Whenever a solution needs exit times, post-order values or ``is $v$ on the current path?'', write the version with \texttt{it[u]} (or recursion).
\end{warningbox}

%------------------------------------------------------------
\subsubsection*{How deep can recursion go?}

Each open frame occupies stack memory, and the stack is small: the operating system reserves a fixed amount for it, independent of the memory limit of the problem. When the recursion goes deeper than the stack allows, the program dies with a \stackoverflow{} --- usually reported as ``runtime error'' with no message at all.

To get real numbers we ran the recursive \texttt{dfs} above on a path $0 - 1 - 2 - \dots - (n-1)$ (depth $= n$) on a Windows machine with MinGW g++ 6.3 and its default 2\,MB stack:

\begin{center}
\small
\begin{tabular}{@{}lccc@{}}
\toprule
\textbf{compiled with} & \textbf{largest $n$ that worked} & \textbf{bytes per frame} & \textbf{$n = 10^6$} \\
\midrule
\texttt{g++ -O2} & about $65\,000$ & $\approx 32$ & crash \\
\texttt{g++ -O0} (no optimisation) & about $43\,000$ & $\approx 48$ & crash \\
\texttt{g++ -O2 -Wl,--stack=268435456} (256\,MB) & --- & $\approx 32$ & works \\
\bottomrule
\end{tabular}
\end{center}

\begin{diagrambox}
\begin{tikzpicture}[x=1cm, y=1cm]
    \draw[->, thick, darkgray] (0,0) -- (12.2,0) node[right, font=\scriptsize] {depth};
    \foreach \x/\l in {0/0, 2/$10^4$, 4/$5{\cdot}10^4$, 6/$10^5$, 8/$2{\cdot}10^5$, 10/$10^6$}{
        \draw[darkgray] (\x,0.08) -- (\x,-0.08); \node[font=\scriptsize, below] at (\x,-0.1) {\l};}
    \fill[successgreen!25] (0,0.25) rectangle (4.6,0.75);
    \node[font=\scriptsize] at (2.3,0.5) {safe everywhere};
    \fill[accentorange!25] (4.6,0.25) rectangle (8.7,0.75);
    \node[font=\scriptsize] at (6.65,0.5) {depends on the judge};
    \fill[warningred!20] (8.7,0.25) rectangle (12,0.75);
    \node[font=\scriptsize] at (10.35,0.5) {only with a huge stack};
    \draw[warningred, thick, dashed] (4.6,0.1) -- (4.6,1.2) node[above, font=\tiny, text=warningred] {2\,MB default: crash near $6.5\cdot10^4$};
\end{tikzpicture}
\figcap{Rule of thumb. Codeforces runs solutions with a large stack (256\,MB), so depth $10^6$ is usually fine there; at home, on many other judges and inside LeetCode's harness the default is a few megabytes.}
\end{diagrambox}

\textbf{What to do.} Before writing a recursive DFS, ask how deep it can go: the worst case is the longest path the DFS can walk, which is $n$ on a path graph. If that can exceed a few $10^4$ and you do not know the judge's stack, use the iterative version. A second, often simpler, way out is the \emph{order array} trick below.

%------------------------------------------------------------
\subsubsection*{Post-order work without recursion: the order array}

Many problems need post-order work only (``compute something for every vertex from its children''). Then a full iterative DFS is not necessary:
\begin{enumerate}[leftmargin=*]
    \item run any iterative search from the root that records \texttt{order} (vertices in the order they are reached) and \texttt{parent};
    \item loop over \texttt{order} \textbf{backwards}: every vertex appears after its parent in \texttt{order}, so going backwards every vertex is processed after all its children, and can add its value to its parent.
\end{enumerate}
\begin{lstlisting}[style=snippet]
for (int i = n - 1; i >= 1; i--) {      // order[0] is the root, it has no parent
    int v = order[i];
    sz[parent[v]] += sz[v];              // all children of v were handled before v
}
\end{lstlisting}
On the DFS forest of $D$ (tree edges $1{\to}2, 2{\to}3, 2{\to}4, 1{\to}5, 5{\to}6, 6{\to}7$; vertex 8 alone) this gives the subtree sizes $\mathrm{sz} = (7,3,1,1,3,2,1,1)$ for vertices $1..8$. The trick works with \emph{any} search order in which parents come before children --- even BFS order --- because it only uses the tree, not the DFS rhythm. It is the workhorse of Pattern~2 when the input is a tree with up to $2\cdot10^5$ vertices.

%------------------------------------------------------------
\subsubsection*{Check your understanding}

\begin{enumerate}[leftmargin=*]
    \item In $D$, change \texttt{adj[1]} to \texttt{6 5 2}. What are the new pre-order and post-order?
    \item While the DFS on $D$ is inside \texttt{dfs(4)}, which frames are open, and which vertices are already left?
    \item Why can variant (b) above hold the same vertex twice on its stack, but the iterator version never?
    \item A tree with $2\cdot 10^5$ vertices is given as a parent array with $p_i < i$. Give an \bigO{n} way to compute all subtree sizes with no search at all.
    \item Your recursive DFS passes the samples and gets ``runtime error'' on test 7, where $n = 2\cdot10^5$. Name the most likely cause and two fixes.
\end{enumerate}

\begin{solutionbox}[Answers]
\begin{enumerate}[leftmargin=*]
    \item DFS from 1 now enters 6, 7, then 5 (6 is done, 5 tries 4), 4, 3, and finally 2 (both its neighbours are done), then 8. Running the program on the changed input prints \textbf{pre: 1 6 7 5 4 3 2 8} and \textbf{post: 7 6 3 4 5 2 1 8}.
    \item Open frames (bottom to top): \texttt{dfs(1)}, \texttt{dfs(2)}, \texttt{dfs(4)}; the path $1 \to 2 \to 4$. Already left: only 3.
    \item Variant (b) pushes every unmarked neighbour, and a vertex can be unmarked when two different vertices push it (6 is pushed by 1 and by 5 in $D$). The iterator version pushes $v$ only at the moment it marks $v$, so each vertex is pushed at most once.
    \item Since $p_i < i$, every child has a larger label than its parent. Loop $i = n, n-1, \dots, 2$ and do \texttt{sz[p[i]] += sz[i]}: the labels themselves are an order with parents before children (the order-array trick with $\texttt{order} = 1,2,\dots,n$).
    \item A stack overflow on a deep tree (a path, or a long ``caterpillar''). Fix 1: the iterative DFS with \texttt{it[u]}. Fix 2: if only post-order work is needed, the order-array trick. (On Codeforces the large stack often hides the problem; on other judges it does not.)
\end{enumerate}
\end{solutionbox}
